\section{Comparison and scope}\label{sec:comparison}
Witsenhausen's static reduction \cite{Witsenhausen} provides an independent reference representation for finite causal experiments. Strategic-measure compactness and information-structure methods \cite{YukselSaldi,YuStrategic} establish powerful existence and optimality principles for larger policy spaces. They do not eliminate the obligation to check that a restriction to one fixed number of \emph{reused} report rows is closed and that its finite responses are continuous. Lemma~\ref{lem:tensor} states that check. Remark~\ref{rem:diagonal} exhibits the failure without the product-history assumption. Once that continuity and a uniform cycle tail have been proved, the attainment and finite-model intersection steps are standard compact analysis, not a new abstract separation principle.

Average-cost partially observed control already has a substantial theory. Yu and Bertsekas \cite{YuBertsekas} study near optimality of the set of finite-state controllers for average-cost POMDPs. Their approximation over finite-state controllers is not a fixed-$M$ robust optimum with a same-$M$ certificate. Y\"uksel's direct history-space approach \cite{YukselHistory} supplies discounted and average-cost existence conditions without first reducing to a belief state, together with finite-window approximation results. Demirci, Kara and Y\"uksel \cite{Demirci} derive average-cost optimality through contraction properties of nonlinear filters. These approaches concern different hypotheses and policy classes; a compulsory observer reset is neither a consequence of filter contraction nor a replacement for their unrestricted-policy conclusions. Our uniform passage in Theorem~\ref{thm:renewal} is instead for a single common stationary finite alphabet, under a return certificate of the physical-plus-observer protocol. Theorem~\ref{thm:geometry} then relates its finite-resource predictions to actual cycle occupation.

The finite-memory approximation results of Kara and Y\"uksel \cite{KaraYuksel} connect finite windows to filter stability. Kim \cite{Kim} gives a policy-conditional Wasserstein comparison under the same closed-loop execution; such a result must not be described as a comparison of independently optimized control policies. This distinction is respected here: Theorem~\ref{thm:renewal} optimizes its full declared common-program class, whereas the geometric prediction Theorem~\ref{thm:geometry} explicitly fixes exploration. Hud\'ak and coauthors \cite{Hudak} extract robust finite-state controllers from learned recurrent policies for hidden-model POMDPs. Their computationally evaluated controllers and the inductive synthesis of Andriushchenko and coauthors \cite{Inductive} address practical synthesis, not the all-Borel coverage lower of Theorem~\ref{thm:compile}. Conversely, our exhaustive bound does not supply their practical scalability. Small-policy complexity \cite{Mundhenk} is a reason to keep finite certifiability separate from efficient synthesis.

The exact renewal theorem admits arbitrary standard Borel hidden dynamics and arbitrary fixed-model families under finite-history product domination and a common return envelope. The effective theorem additionally needs computable primitive or response certificates; the displayed hidden and quantum implications supply them for the stated subclasses. The geometric theorem needs actual mass, a global cover, and candidate-valid occupation stability. The direct holding theorem covers nondominated preparation families but does not remove domination from general adaptive control. These are distinct levels of assumptions, not interchangeable uses of the word ``general.''

The remaining mathematical problems include an unrestricted nonreset average-risk theorem under comparably primitive hypotheses; a broad nondominated adaptive common-controller theory; a near-necessary intrinsic criterion replacing global covers and sufficient stability; efficient synthesis; and a complete matched region for program length, scratch, precision, planning time, physical time and minimal simulator state. No theorem here asserts a large-deviation principle from total-variation proximity, an infinite-path law from finite domination, or an unbounded-operator closure from finite matrices. The earlier discounted theorem \cite{Discounted} and earlier exact-deadline and causal-allocation results retain their own protocols. They are preserved as complete technical companions, not used as unproved premises of the new average-risk theorem.
